\section{Introduction}
\label{sec:intro}

Science has always relied on a few prestigious venues to certify that a result is worth reading. Hiring committees, funders, and individual scholars all use the venue label as a coarse signal: a paper at \iclr or in \emph{Nature} is presumed to clear a quality bar that an arbitrary preprint does not. This signaling function rests on a simple supply-demand premise---a fixed number of slots, a much larger pool of submissions, and a screening mechanism (peer review) that is reliable enough to separate the two.

Submission volume has grown dramatically. Modern AI conferences now exceed $10{,}000$ submissions per venue per year~\citep{yang2025papercopilot}; \iclr alone grew almost $300\times$ from $67$ submissions in $2013$ to $19{,}814$ in $2026$. Across science as a whole, output doubles every $8$--$13$ years~\citep{bornmann2015growth}. As volume grows, three things might break the signaling premise: the gatekeeper may no longer have enough capacity to be informative; reviewers, asked to handle more papers each, may grow noisier; and even when per-paper precision is preserved, downstream consumers (readers and hiring committees) may no longer be able to use the signal because their own budgets have not grown. {\bf Our main question is: how rapidly does the signaling effect of gate-keeping diminish with $\Nsub$, and through what mechanisms?}

\para{Why this matters.} If volume growth silently destroys the certification function, three things break. (i) Readers cannot allocate attention---knowing a paper is at \iclr no longer narrows the choice from $19{,}814$ to a manageable list. (ii) Hiring and tenure committees cannot identify talent---publication count in elite venues becomes uninformative about underlying ability. (iii) The system as a whole cannot reward quality, weakening the incentives that hold scientific careers together. \citet{heckman2019tyranny} already document that the \emph{Top-Five} screen in economics is far from reliable; we ask whether the situation is systematically worsening as the field grows.

\para{Gap.} Existing work touches all three pieces but does not connect them. \citet{akerlof1970lemons}, \citet{spence1973job}, and \citet{long2023lemons} give static signaling models that do not vary submission volume $\Nsub$ as the explanatory variable. \citet{heckman2019tyranny} and \citet{kwiek2026topperformers} describe gatekeeping at single snapshots. \citet{bornmann2015growth} and \citet{fu2021disproportion} quantify volume growth and per-paper information drop without a signaling mechanism. Agent-based simulators of peer review such as \textsc{AgentReview}~\citep{jin2024agentreview} and \citet{grimaldo2018reputation} fix load. {\bf No prior work combines capacity-constrained Bayesian signaling, a reader consumption model, and a researcher-level identifiability analysis into a single framework parameterized by $\Nsub$.}

\para{Our approach.} We instantiate the gate-keeping system at three layers. At the {\bf paper layer}, $\Nsub$ papers with latent qualities $\quality_i$ produce noisy reviewer scores $\score_i = \quality_i + \varepsilon_i$ with $\varepsilon_i \sim \gN(0,\sigrev^2)$, and the gatekeeper accepts the top $\Kcap$. At the {\bf reader layer}, $M$ readers with fixed budget $\budget$ papers per year choose what to read from the venue's accepted pool. At the {\bf researcher layer}, $M_r$ researchers with latent ability $\ability_i$ submit $S$ papers per year for $T$ years, and we measure how well their publication count ranks $\ability_i$. We derive analytical and numerical results across distributions, capacity policies, and reviewer-noise regimes; we then calibrate the model on $14$ years of \iclr submission and review data extracted via \papercopilot~\citep{yang2025papercopilot}.

\para{Quantitative preview.} We find three compounding decay mechanisms. (1) Under fixed capacity, mutual information $\miquant$ and recall on the true top-1\% fall monotonically in $\Nsub$ (Kendall $\tau \le -0.95$, $p<10^{-6}$ for every $(\sigrev, F)$ cell). (2) The Spearman correlation between researcher ability and publication count collapses from $0.78$ at $M{=}200$ researchers to $0.26$ at $M{=}20{,}000$, with $95.5\%$ of researchers publishing nothing. (3) On \iclr 2017--2026, the standardized posterior contrast $\hcontrast / \stdrating$ falls log-linearly in $\Nsub$ with slope $-0.32$ per $e$-fold ($R^2 = 0.85$, $p = 1.1\times10^{-3}$), and the mean per-paper reviewer SD rises from $0.86$ to $1.53$. The model recovers \citet{akerlof1970lemons}, \citet{spence1973job}, and \citet{long2023lemons} in the no-capacity-constraint limit and matches \citet{heckman2019tyranny}'s observation that the top-tier screen captures only about half of the truly influential work.

\para{Contributions.}
\begin{itemize}[leftmargin=*,itemsep=0pt,topsep=0pt]
    \item We propose a three-layer (paper / reader / researcher) capacity-constrained Bayesian signaling model parameterized by submission volume $\Nsub$, and we derive its decay behavior in closed form for Gaussian quality and numerically for log-normal and Pareto quality.
    \item We show that the choice of capacity policy ($\Kcap$ constant, $\Kcap \propto \sqrt{\Nsub}$, or $\Kcap \propto \Nsub$) is the central design lever, and that no policy simultaneously preserves per-paper precision, researcher identifiability, and reader-level relevance.
    \item We extract submission counts and per-paper reviewer-rating dispersion from $14$ years of \iclr data and show that the predicted log-linear decay of the standardized posterior contrast holds with $R^2 = 0.85$.
    \item We release reproducible code (analytical sweeps, Monte Carlo simulators, empirical extractors) that runs end-to-end on CPU in under one minute.
\end{itemize}

The paper proceeds as follows. \secref{sec:related} positions the work; \secref{sec:method} formalizes the three-layer model and the empirical pipeline; \secref{sec:results} reports analytical, simulation, and empirical findings; \secref{sec:discussion} discusses limitations and policy implications; \secref{sec:conclusion} concludes.
